\documentclass[11pt]{article}
\usepackage[margin=1in]{geometry}
\usepackage{amsmath,amssymb,amsthm,booktabs,array}
\usepackage[T1]{fontenc}
\usepackage{lmodern,microtype}
\usepackage{xurl}
\usepackage[colorlinks=true,linkcolor=blue,citecolor=blue,urlcolor=blue]{hyperref}
\newtheorem{theorem}{Theorem}[section]
\newtheorem{proposition}[theorem]{Proposition}
\newtheorem{lemma}[theorem]{Lemma}
\newtheorem{corollary}[theorem]{Corollary}
\theoremstyle{remark}
\newtheorem{remark}[theorem]{Remark}
\newcommand{\dd}{\,d}
\newcommand{\R}{\mathbb R}
\newcommand{\C}{\mathbb C}
\newcommand{\Vcal}{\mathcal V_g}
\newcommand{\Rea}{\operatorname{Re}}
\newcommand{\Ima}{\operatorname{Im}}
\newcommand{\norm}[1]{\left\|#1\right\|}
\DeclareMathOperator{\supp}{supp}
\numberwithin{equation}{section}
\title{Prime response variance exponents and the Riemann hypothesis}
\author{Edward Baker}
\date{Draft, 4 October 2026}
\begin{document}
\maketitle
\begin{abstract}
We fix an explicit, compactly supported, prepared prime response and study
its variance on complete dyadic intervals. For each $0\le\delta\le1$, a
global bound $O(X^{2+\delta})$ is equivalent to every nontrivial zeta zero
lying in the closed strip $|\Rea\rho-1/2|\le\delta/2$. The proof uses an
elementary noncancellation argument and the true Laplace transform of the
response. Thus bounds along any positive sequence of exponents tending
to zero are equivalent to the Riemann hypothesis, without uniformity of
their constants. We give an unconditional baseline, transfer existing
finite certificates to quadratic bounds through $10^{102}$, and reduce
the global problem to an exact signed short-interval covariance with
quadratic approximation cost. A finite-prefix construction shows why
preparation, prime-number-theorem estimates, and diagonal controls alone
do not imply an exponent descent. No fixed global exponent below one,
new zero-free strip, or actual-prime descent rule is established.
\end{abstract}

\section{The fixed response and the exponent program}
Let $\Lambda$ be the von Mangoldt function and
$\psi(t)=\sum_{2\le n\le t}\Lambda(n)$, with $\psi(t)=0$ for $t<2$.
All logarithms are natural. Fix, once and for all,
\begin{equation}\label{eq:probe}
\begin{gathered}
a=\tfrac14,\qquad h(v)=(1-16v^2)^8\mathbf1_{|v|<a},\qquad
g_0=-h'''+\tfrac14h',\qquad g=g_0/\sqrt\nu,\\
\nu=\norm{g_0}_2^2=\frac{146640624550936576}{37921101075},\qquad
A=e^{-a},\quad B=e^a,\qquad w(t)=t^{-1/2}g(-\log t)\quad(t>0).
\end{gathered}
\end{equation}
The functions $h,g,w$ are extended by zero off their indicated supports;
in particular $\supp w=[A,B]$. The preparation will imply
$\int w=0$ and $\int w^2=1$. Define
\begin{equation}\label{eq:response}
\begin{split}
p_g(y)&=\sum_{n\ge2}\frac{\Lambda(n)}{\sqrt n}g(y-\log n),\\
V_g(x)&=\sum_{n\ge2}\Lambda(n)w(n/x)=\sqrt x\,p_g(\log x),\qquad
\Vcal(X)=\int_X^{2X}|V_g(x)|^2\dd x.
\end{split}
\end{equation}
These sums are locally finite. For $x\in[X,2X]$ every active integer
lies in $[AX,2BX]$. The entire support window is always retained;
weights vanish at its endpoints, so endpoint inclusion is immaterial.
The word variance denotes this centered quadratic integral, without
division by the interval length.

\begin{theorem}[Fixed exponent and closed zero strip]\label{thm:main}
For each fixed $\delta\in[0,1]$, the following are equivalent:
\begin{align}
\Vcal(X)&=O(X^{2+\delta})\quad\text{for all sufficiently large real }X;
\label{eq:global}\\
|\Rea\rho-\tfrac12|&\le\tfrac\delta2
\quad\text{for every nontrivial zero $\rho$ of $\zeta$}.
\label{eq:strip}
\end{align}
The constants and starting threshold in \eqref{eq:global} may depend on
$\delta$. Zeros on the boundary of the strip are permitted.
\end{theorem}

\begin{corollary}\label{cor:rh}
The exponent $\delta=1$ is unconditional, and $\delta=0$ is equivalent
to the Riemann hypothesis (RH). Each of the following is also equivalent
to RH: \eqref{eq:global} holds for every $\delta\in(0,1]$; or it holds
for some positive sequence $\delta_j\to0$. No uniform constants or
uniform starting thresholds are required. A single fixed $\delta<1$
already gives an all-height zero-free strip; for example, $\delta=0.99$
would exclude $\Rea\rho>0.995$.
\end{corollary}

The illustrative exponent $0.99$ has not been attained here. The first
global milestone is any fixed exponent below one for the actual
coefficients $\Lambda(n)$. A rule improving an already proved exponent
is a separate open problem. Even an infinite sequence of improvements
with positive limiting exponent gives only the corresponding fixed
strip. A finite range or a fitted numerical slope does not meet the
global quantifier in \eqref{eq:global}.

For completeness, the optimal exponent is
\begin{equation}\label{eq:optimal}
\delta_*:=\inf\{\delta\in[0,1]:\Vcal(X)=O(X^{2+\delta})\}
=2\sup_\rho(\Rea\rho-\tfrac12).
\end{equation}
This translates Theorem~\ref{thm:main}; it does not compute the supremum.
Indeed the supremum itself defines a closed strip, so the theorem also
supplies the variance bound at this optimal exponent.

The purpose is a focused quantitative reformulation and research
program. Sections~\ref{sec:probe}--\ref{sec:equivalence} prove the
equivalence directly. Section~\ref{sec:bounds} separates global
baseline estimates from finite quadratic certificates.
Section~\ref{sec:arithmetic} gives the principal arithmetic formulation,
and Section~\ref{sec:milestone} identifies the missing global input.
The fixed probe and certificate inputs arose in the parent
Sonin critical-boundary investigation; no operator results from that
investigation are needed for the proofs below. No claim of mathematical
priority is made.

\section{Preparation, noncancellation, and transform decay}\label{sec:probe}
Use exclusively the plus-exponent convention
\begin{equation}\label{eq:transforms}
G(z)=\int_{-a}^a g(v)e^{zv}\dd v,\qquad
H(z)=\int_{-a}^a h(v)e^{zv}\dd v.
\end{equation}
Both transforms are entire. In sources using a minus-exponent transform,
that transform is $G(-z)$ in the present notation.

\begin{lemma}[Probe identities]\label{lem:probe}
The zero extension of $g$ is real, odd, $C^4$, and normalized in $L^2$.
It and its derivatives through order four vanish at $\pm a$. Moreover,
\begin{equation}\label{eq:moments}
G(0)=G(\tfrac12)=G(-\tfrac12)=0,\qquad
\int_\R w(t)\dd t=0,\qquad \int_\R w(t)^2\dd t=1,
\end{equation}
and
\begin{equation}\label{eq:series}
G(z)=\frac{z(z^2-1/4)}{\sqrt\nu}H(z),\qquad
\frac{H(z)}{H(0)}=
\sum_{k=0}^\infty\frac{(z^2/64)^k}{(19/2)_k\,k!}.
\end{equation}
Here $(q)_k=q(q+1)\cdots(q+k-1)$ and $(q)_0=1$.
\end{lemma}
\begin{proof}
The polynomial $h$ has a zero of order eight at each endpoint; its zero
extension is $C^7$. Its third derivative therefore vanishes to order five
there, giving the stated regularity of $g$. Parity and reality follow
from the definition. Exact integration of the polynomial square gives
$\nu$ in \eqref{eq:probe}; an explicit finite-arithmetic prescription is
given in Appendix~\ref{app:certificate}.

Three integrations by parts, with zero boundary terms, give the first
identity in \eqref{eq:series} and hence the three vanishing moments.
Expanding $e^{zv}$ and integrating the even monomials proves the second
identity: substitution $u=4v$, followed by $r=u^2$, gives
\[
\frac{\int_{-a}^a v^{2k}h(v)\dd v}{\int_{-a}^a h(v)\dd v}
=16^{-k}\frac{(1/2)_k}{(19/2)_k},\qquad
(2k)!=4^k k!(1/2)_k.
\]
Finally $t=e^{-v}$ gives $\int w=G(-1/2)$ and
$\int w^2=\int g^2=1$. In particular $w\in C_c^4(\R)$.
\end{proof}

\begin{lemma}[Elementary noncancellation]\label{lem:noncancel}
Every zero of $H$ is nonzero and purely imaginary. Consequently
$G(-s)\ne0$ whenever $0<\Rea s<1/2$. Its zero at $s=1/2$ is simple.
\end{lemma}
\begin{proof}
Put $f=H/H(0)$. If $f(z_0)=0$, then $z_0\ne0$. For
$u(t)=f(z_0t)$, the series \eqref{eq:series} gives
\[
(t^{18}u')'=\frac{z_0^2}{16}t^{18}u,\qquad
u(0)=1,\quad u'(0)=0,\quad u(1)=0.
\]
Multiply by $\overline u$ and integrate over $[0,1]$:
\[
-\int_0^1t^{18}|u'|^2\dd t
=\frac{z_0^2}{16}\int_0^1t^{18}|u|^2\dd t.
\]
The boundary terms vanish at both ends. Both integrals are strictly
positive before the minus sign, since $u$ is neither zero nor constant.
Thus $z_0^2$ is strictly negative real. The polynomial factor in
\eqref{eq:series} proves noncancellation. Since $H(-1/2)>0$,
the preparation zero at $s=1/2$ is simple.
\end{proof}

\begin{lemma}[Sixth-power decay]\label{lem:decay}
There is a fixed $K>0$ such that
\begin{equation}\label{eq:decay}
|G(z)|\le K e^{a|\Rea z|}|z|^{-6}\quad(z\ne0).
\end{equation}
In particular, $G(\sigma+it)=O_g((1+|t|)^{-6})$ uniformly for
$|\sigma|\le1/2$.
\end{lemma}
\begin{proof}
The piecewise fifth derivative of $g$ has finite endpoint jumps.
Its sixth distributional derivative is a finite signed measure,
consisting of an interior polynomial density and two endpoint atoms.
Distributional integration by parts gives
$z^6G(z)=\int e^{zv}\dd(D^6g)(v)$.
Take $K=\norm{D^6g}_{\rm TV}$, or the larger explicit constant in
Appendix~\ref{app:certificate}. Boundedness on compact sets covers
small $t$ in the final assertion.
\end{proof}

\section{Proof of the exponent equivalence}\label{sec:equivalence}
We use the classical meromorphic continuation and functional equation of
$\zeta$, its simple pole at one, its simple trivial zeros $-2k$,
and the Dirichlet series
\begin{equation}\label{eq:dirichlet}
-\frac{\zeta'}{\zeta}(z)=\sum_{n\ge2}\Lambda(n)n^{-z}
\qquad(\Rea z>1).
\end{equation}
Nontrivial zeros satisfy $0<\Rea\rho<1$ and are invariant, with
multiplicity, under conjugation and $\rho\mapsto1-\rho$.
Write $N(T)$ for their number with $0<\Ima\rho\le T$, counted with
multiplicity. The classical zero count gives
$N(T)=O(T\log(T+2))$ and $N(T+1)-N(T)=O(\log(T+2))$.
For contour shifts one may choose heights $T_j\to\infty$ avoiding zero
ordinates such that $\zeta'/\zeta(\sigma\pm iT_j)=O(\log^2T_j)$ on
each fixed bounded $\sigma$ interval. This follows from the local count
and the logarithmic-derivative partial fraction formula, by staying
at distance $\gg1/\log T_j$ from the ordinates. These are standard
zeta inputs; see \cite{DLMF,ExplicitFormula,HSW}.
All sums indexed by $\rho$ below run over distinct nontrivial zeros,
include both ordinate signs, and carry their multiplicities $m_\rho$.

\begin{proposition}[Complete linear explicit formula]\label{prop:explicit}
For $y>a$,
\begin{equation}\label{eq:explicit}
p_g(y)=-\sum_\rho m_\rho G(1/2-\rho)e^{(\rho-1/2)y}
-\sum_{k\ge1}G(2k+1/2)e^{-(2k+1/2)y}.
\end{equation}
Both series converge absolutely and locally uniformly on $y>a$.
The nontrivial coefficients alone are absolutely summable.
\end{proposition}
\begin{proof}
Fourier inversion for the compact $C^4$ function $g$, on the line
$\Rea z=2$, and absolute convergence of \eqref{eq:dirichlet} give
\[
p_g(y)=\frac1{2\pi i}\int_{2-i\infty}^{2+i\infty}
-\frac{\zeta'}{\zeta}(z)
e^{(z-1/2)y}G(1/2-z)\dd z.
\]
The integrals and the interchange with the Dirichlet series converge
absolutely by \eqref{eq:decay}. Shift to
$\Rea z=-(2M+1)$, first sending the horizontal sides to the good heights
just described. On a fixed strip their integrands are
$O_{M,y}(T_j^{-6}\log^2T_j)$, so the horizontal integrals vanish.
The pole at one would contribute $e^{y/2}G(-1/2)=0$.
At a nontrivial zero the residue is
$-m_\rho e^{(\rho-1/2)y}G(1/2-\rho)$, and at $-2k$ it is
$-e^{-(2k+1/2)y}G(2k+1/2)$.

On the negative odd vertical lines, the functional equation and the
Dirichlet series in $\Rea(1-z)>1$ give
$\zeta'/\zeta(z)=O(\log(|z|+2))$, uniformly as $M$ increases.
The exponential in the remaining integral is bounded by
$e^{-(2M+3/2)(y-a)}$ after applying \eqref{eq:decay}; the remaining
integral of $|z-1/2|^{-6}\log(|z|+2)$ is
$O(M^{-5}\log(M+2))$. Thus this vertical integral tends to zero for
$y>a$. This contour argument works directly for the present $C^4$
probe and needs no unstated smoothness assumption.

The zero count and \eqref{eq:decay} imply
$\sum_\rho m_\rho|G(1/2-\rho)|<\infty$.
The factors $e^{(\Rea\rho-1/2)y}$ are uniformly bounded on compact
$y$ intervals. Finally $\norm g_1\le\sqrt{2a}$ gives
\begin{equation}\label{eq:trivial}
\sum_{k\ge1}|G(2k+1/2)|e^{-(2k+1/2)y}
\le \frac{\sqrt{1/2}\,e^{-5(y-a)/2}}{1-e^{-2(y-a)}}\quad(y>a).
\end{equation}
This proves the convergence assertions, including local uniformity.
\end{proof}

\begin{proof}[Proof of Theorem~\ref{thm:main}]
Assume \eqref{eq:global}. The exact change of variables $x=e^y$ gives
\begin{equation}\label{eq:shell}
\int_{\log X}^{\log(2X)}|p_g(y)|^2\dd y
=\int_X^{2X}\frac{|V_g(x)|^2}{x^2}\dd x
\le X^{-2}\Vcal(X).
\end{equation}
Let $\ell=\log2$. On the shells $[j\ell,(j+1)\ell]$ this is
$O(2^{j\delta})$ for large $j$. For $\sigma=\Rea s>\delta/2$,
Cauchy--Schwarz therefore gives
\[
\int_{j\ell}^{(j+1)\ell}e^{-\sigma y}|p_g(y)|\dd y
\le \ell^{1/2}2^{-j\sigma}
\left(\int_{j\ell}^{(j+1)\ell}|p_g(y)|^2\dd y\right)^{1/2}
\ll 2^{j(\delta/2-\sigma)}.
\]
The sum is geometric, locally uniformly in $\Rea s>\delta/2$.
On each finite interval $p_g$ is continuous by local finiteness.
Consequently the true integral
\begin{equation}\label{eq:laplace}
P_g(s)=\int_0^\infty e^{-sy}p_g(y)\dd y
\end{equation}
is absolutely convergent and holomorphic in that half-plane; local
uniform domination, or Morera's theorem, justifies holomorphy.

Since $a<\log2$, $p_g(y)=0$ for $0\le y<\log2-a$, and every translated
packet has its entire support above zero. For $\Rea s>1/2$,
absolute termwise integration thus gives exactly
\begin{equation}\label{eq:laplace-zeta}
P_g(s)=G(-s)\sum_{n\ge2}\Lambda(n)n^{-1/2-s}
=-G(-s)\frac{\zeta'}{\zeta}(1/2+s).
\end{equation}
There is no initial endpoint correction. Uniqueness of meromorphic
continuation identifies the right side with the holomorphic function
\eqref{eq:laplace} throughout $\Rea s>\delta/2$.
The pole of $\zeta$ at $s=1/2$ is canceled by preparation.
If $\rho=\beta+i\gamma$ had $\beta-1/2>\delta/2$, then
$0<\Rea(\rho-1/2)<1/2$ and Lemma~\ref{lem:noncancel} would leave the
nonzero residue $-m_\rho G(1/2-\rho)$ at $s=\rho-1/2$.
This contradicts holomorphy. Functional-equation symmetry supplies the
left boundary, proving the closed strip \eqref{eq:strip}.

Conversely, assume that strip. Absolute summability in
Proposition~\ref{prop:explicit} bounds the nontrivial sum by
$C_g e^{\delta y/2}$ for $y\ge1$. The trivial sum is bounded there
by \eqref{eq:trivial}. Hence
$p_g(y)=O_g(e^{\delta y/2})$, and
$V_g(x)=O_g(x^{(1+\delta)/2})$. Integrating its square over $[X,2X]$
proves \eqref{eq:global}. The initial compact interval is controlled
by the original finite arithmetic sum.

Both arguments include $\delta=0$ and $\delta=1$. The converse uses
a non-strict bound on each real part and therefore permits boundary
zeros, with no logarithmic loss. The forward argument only uses
$\Rea s>\delta/2$; no convergence at the boundary is asserted.
\end{proof}

\begin{proof}[Proof of Corollary~\ref{cor:rh}]
Every nontrivial zero lies in $0<\Rea\rho<1$, giving $\delta=1$.
At $\delta=0$ the strip is exactly the critical line. The bounds along
$\delta_j\to0$ imply all the corresponding closed strips; their
intersection is the critical line. Conversely RH gives the quadratic
bound and hence all larger exponents. Constants may vary with $j$
because the intersection is a statement about each fixed zero.
\end{proof}

\begin{remark}
If $J(Y)=\int_0^Y|p_g(y)|^2\dd y$ is also considered, summing
\eqref{eq:shell} gives $J(Y)=O(1+Y)$ at $\delta=0$, not $O(1)$.
For positive $\delta$ it gives $J(Y)=O(e^{\delta Y})$.
Neither observation is needed to assert a transform on its boundary.
\end{remark}

\section{Global baseline and finite quadratic bounds}\label{sec:bounds}
\subsection{Transfer of a prime-number-theorem error}
Put $E(t)=\psi(t)-t$ and $M_w=\int_A^B u|w'(u)|\dd u$.
Stieltjes integration by parts and $\int w=0$ give
\begin{equation}\label{eq:pnt-transfer}
V_g(x)=\int w(t/x)\dd(\psi(t)-t)
=-\int_A^B E(xu)w'(u)\dd u.
\end{equation}
The weight vanishes at both endpoints; the identity retains the entire
support even when an endpoint is an integer.

\begin{proposition}[Unconditional envelope transfer]\label{prop:baseline}
If $|E(t)|\le t r(t)$ throughout $[AX,2BX]$, where $r\ge0$, then
\begin{equation}\label{eq:baseline}
\Vcal(X)\le\frac73M_w^2X^3
\left(\sup_{AX\le t\le2BX}r(t)\right)^2.
\end{equation}
In particular, unconditionally as $X\to\infty$,
\begin{equation}\label{eq:jy}
\Vcal(X)=O_g\left(X^3(\log X)^{3.602}
\exp\left[-0.3706\frac{(\log X)^{3/5}}{(\log\log X)^{1/5}}\right]\right).
\end{equation}
\end{proposition}
\begin{proof}
Equation~\eqref{eq:pnt-transfer} bounds $|V_g(x)|$ by
$xM_w\sup r$ on the shell. Integrating $x^2$ gives $7X^3/3$.
Johnston--Yang \cite[Theorem 1.4]{JY} prove, for $t\ge23$,
\[
|E(t)|\le t\,0.026(\log t)^{1.801}
\exp\left[-0.1853\frac{(\log t)^{3/5}}{(\log\log t)^{1/5}}\right].
\]
For $X\ge23/A$, inserting this exact function $r$ into
\eqref{eq:baseline} is an explicit bound with the displayed supremum.
For the asymptotic statement, put
$F(X)=(\log X)^{3/5}/(\log\log X)^{1/5}$.
The mean value theorem in $\log X$ gives
$F(cX)-F(X)=o(1)$ uniformly for $c\in[A,2B]$; the logarithmic
prefactor ratio also tends uniformly to one. Squaring the envelope
proves \eqref{eq:jy}.
\end{proof}

No monotonicity of the explicit envelope is needed. The decaying factor
in \eqref{eq:jy} is subpower: $F(X)=o(\log X)$. It supplies no fixed
improvement of the global exponent $\delta=1$. This is an application
of an existing prime-number-theorem estimate, not a new arithmetic
saving. It does not assert that the actual variance attains the upper
comparison function.

\subsection{Finite verification and the linear tail}
The following finite conclusions use published critical-line verification
and retained outward numerical records. Their analytic transfer is given
here; Appendix~\ref{app:certificate} identifies the records and exact
rational checks.

\begin{theorem}[Finite continuous quadratic bounds]\label{thm:finite}
For every real $X$ in the indicated intervals,
\begin{equation}\label{eq:finite}
\begin{aligned}
\Vcal(X)&<38X^2 &&(e\le X\le10^{99}),\\
\Vcal(X)&<40X^2 &&(e\le X\le10^{100}),\\
\Vcal(X)&<72X^2 &&(e\le X\le10^{102}).
\end{aligned}
\end{equation}
These bounds require no global RH assumption.
\end{theorem}
\begin{proof}
Take $L=500$, $H=3\cdot10^{12}$, and
$K=9470610144359/\sqrt\nu$ in \eqref{eq:decay}.
Hasanalizade--Shen--Wong \cite[Corollary 1.2]{HSW} imply
$N(t)\le t\log t$ for $t\ge100$; the explicit comparison is recorded
in Appendix~\ref{app:certificate}. Stieltjes partial summation gives
\begin{equation}\label{eq:tail}
\sum_{|\Ima\rho|>T}m_\rho|\Ima\rho|^{-6}
\le12\int_T^\infty\frac{\log t}{t^6}\dd t
=12T^{-5}\left(\frac{\log T}{5}+\frac1{25}\right),\quad T\ge100.
\end{equation}
Indeed for positive ordinates the exact boundary term is
$-T^{-6}N(T)\le0$; the factor two includes both signs.

Platt--Trudgian \cite[Theorem 1]{PT} verify the critical-line location
through a height exceeding $H$. The retained records certify
$N(500)=269$ and
\[
S_{500}=2\sum_{j=1}^{269}|G(-i\gamma_j)|
=4.813490676290039592\ldots,
\]
with outward enclosures rather than merely decimal approximations.
For the critical zeros between $L$ and $H$, apply \eqref{eq:tail}
with $T=L$ and no exponential real-part factor. For the zeros above
$H$, use $|\Rea\rho-1/2|<1/2$ and
$|e^{(\rho-1/2)y}|\le e^{y/2}$ for $y\ge1$.
The trivial zeros are bounded by \eqref{eq:trivial}. Thus, for every
$y\ge1$ without an upper restriction,
\begin{equation}\label{eq:envelope}
|p_g(y)|\le b+R_H e^{y/2},
\end{equation}
where
\begin{equation}\label{eq:b-r}
\begin{split}
b&=S_{500}+12K L^{-5}\left(\frac{\log L}{5}+\frac1{25}\right)
+\frac{\sqrt{1/2}\,e^{-15/8}}{1-e^{-3/2}},\\
R_H&=12Ke^{1/8}H^{-5}\left(\frac{\log H}{5}+\frac1{25}\right).
\end{split}
\end{equation}
The infinite counting majorant used for the finite segment $(L,H]$
does not assume criticality beyond $H$.
The outward inputs prove the strict rational bounds
\begin{equation}\label{eq:rounding}
b<\frac{124}{25}=4.96,\qquad
R_H<\frac{39}{25}\,10^{-51}=1.56\cdot10^{-51}.
\end{equation}
Since $|V_g(x)|\le b\sqrt{x}+R_Hx$ for $x\ge e$, exact integration
over the full shell gives
\begin{equation}\label{eq:budget}
\Vcal(X)\le\frac32b^2X^2+
\frac45(2^{5/2}-1)bR_HX^{5/2}+\frac73R_H^2X^3.
\end{equation}
After division by $X^2$, this positive budget increases with $X$.
Insert \eqref{eq:rounding}, use
$\frac45(2^{5/2}-1)<466/125$, and evaluate the rational upper budgets
at the three right endpoints. They are respectively
$37.8311431296$, $39.84376128$, and $71.4265728$.
Appendix~\ref{app:certificate} gives the exact fractions.
Monotonicity proves all three continuous ranges, including endpoints.
\end{proof}

\begin{corollary}[A systematic finite-range milestone]\label{cor:height}
Keep the low cutoff and the constant $b$ fixed. For a ceiling
$C>3b^2/2$, let $q_C>0$ solve
\[
\tfrac73q_C^2+\tfrac45(2^{5/2}-1)bq_C+\tfrac32b^2=C.
\]
Along heights $H\ge500$ through which criticality is rigorously verified,
\eqref{eq:budget} proves $\Vcal(X)\le CX^2$ for
$e\le X\le(q_C/R_H)^2$. The resulting range satisfies
\begin{equation}\label{eq:height-law}
X_{\max}\asymp_{g,C} H^{10}/(\log H)^2.
\end{equation}
\end{corollary}
\begin{proof}
Substitute $q=R_H\sqrt X$ in the normalized budget. Formula
\eqref{eq:b-r} gives $R_H\asymp_g H^{-5}\log H$.
\end{proof}
The linear coefficient tail is of order $H^{-5}\log H$; a
twelfth-power transform bound for an autocorrelation cannot be used
for this linear response. At any fixed $H$, \eqref{eq:budget} retains
a cubic term as $X\to\infty$. Therefore \eqref{eq:height-law} improves
a finite range, and does not establish a global quadratic bound.

\section{Short intervals and an exact signed covariance}\label{sec:arithmetic}
Fix a sufficiently large real $X$ and $1\le h\le AX$. Define
\begin{equation}\label{eq:short-defs}
\begin{split}
D_h(t)&=\psi(t+h/2)-\psi(t-h/2)-h,\\
S(X,h)&=\int_{AX}^{2BX}|D_h(t)|^2\dd t,\\
V_h(x)&=\frac1h\int_\R w(t/x)D_h(t)\dd t,\qquad
Q_h(X)=\int_X^{2X}|V_h(x)|^2\dd x.
\end{split}
\end{equation}
The centered difference uses the actual prime-power counting function.
All integrals in a fixed window are finite.

\begin{proposition}[Complete-support transfer]\label{prop:short}
Set
\begin{equation}\label{eq:capp}
C_{\rm app}=\frac{(4\log2)(B+A/2)}{24}\norm{w''}_\infty.
\end{equation}
Uniformly in the indicated real $X,h$ range,
\begin{align}
\Vcal(X)&\le\frac{3X^2}{h^2}S(X,h)+C_{\rm app}^2\frac{h^4}{X},
\label{eq:gate}\\
\left|\sqrt{\Vcal(X)}-\sqrt{Q_h(X)}\right|
&\le C_{\rm app}\frac{h^2}{\sqrt{2X}}.
\label{eq:norm-error}
\end{align}
Moreover the exact signed covariance is
\begin{equation}\label{eq:covariance}
\begin{split}
Q_h(X)&=\frac1{h^2}\int_{AX}^{2BX}\int_{AX}^{2BX}
D_h(t)D_h(u)W_X(t,u)\dd t\dd u,\\
W_X(t,u)&=\int_X^{2X}w(t/x)w(u/x)\dd x.
\end{split}
\end{equation}
For every fixed $\delta\ge0$,
\begin{equation}\label{eq:Q-equivalence}
Q_{X^{3/4}}(X)=O(X^{2+\delta})
\quad\Longleftrightarrow\quad\Vcal(X)=O(X^{2+\delta}).
\end{equation}
\end{proposition}
\begin{proof}
First, finite Fubini and exact centering give
\begin{equation}\label{eq:finite-fubini}
V_h(x)=\sum_{n\ge2}\Lambda(n)\overline w_{x,h}(n),\qquad
\overline w_{x,h}(n)=\frac1h\int_{-h/2}^{h/2}w((n+u)/x)\dd u.
\end{equation}
For a fixed $n$, the inverse set of $t$ for the counting interval
$(t-h/2,t+h/2]$ is $[n-h/2,n+h/2)$, with endpoints irrelevant to
Lebesgue integration. The density term vanishes because
$\int w(t/x)\dd t=x\int w=0$. This proves the identity for
noninteger $X,h$ without a rounding error.

Taylor's theorem for the $C^2$ zero extension, followed by symmetric
averaging, gives for every real $n$
\begin{equation}\label{eq:taylor}
|\overline w_{x,h}(n)-w(n/x)|
\le \frac{h^2}{24x^2}\norm{w''}_\infty.
\end{equation}
The difference vanishes off $[Ax-h/2,Bx+h/2]$. Thus the estimate also
includes every extra atom created near the original support endpoints.
For clarity, the elementary Chebyshev constant used here is
\begin{equation}\label{eq:cheb}
\psi(T)\le(4\log2)T\qquad(T\ge1).
\end{equation}
Indeed Legendre's valuation formula implies
$\psi(2n)-\psi(n)\le\log\binom{2n}{n}\le2n\log2$.
Sum at dyadic integer values and round $T$ upward to the next power
of two to obtain \eqref{eq:cheb}.

Since $h\le AX\le Ax$, \eqref{eq:taylor} and \eqref{eq:cheb} yield
\[
|V_h(x)-V_g(x)|\le\frac{h^2\norm{w''}_\infty}{24x^2}
\psi(Bx+h/2)\le C_{\rm app}\frac{h^2}{x}.
\]
Integration gives
\begin{equation}\label{eq:L2-error}
\int_X^{2X}|V_h-V_g|^2\dd x\le C_{\rm app}^2\frac{h^4}{2X}.
\end{equation}
On the other hand, $\int w(t/x)^2\dd t=x$ and support in $[Ax,Bx]$
give by Cauchy--Schwarz
\[
|V_h(x)|^2\le\frac{x}{h^2}\int_{Ax}^{Bx}|D_h(t)|^2\dd t
\le\frac{x}{h^2}S(X,h),\qquad
Q_h(X)\le\frac{3X^2}{2h^2}S(X,h).
\]
Apply $|r+s|^2\le2|r|^2+2|s|^2$ to prove \eqref{eq:gate}, and
the reverse triangle inequality in $L^2([X,2X])$ to prove
\eqref{eq:norm-error}. Finite Fubini gives \eqref{eq:covariance},
since the union of the $t$ supports is exactly $[AX,2BX]$.
Finally $h=X^{3/4}$ is admissible for all sufficiently large $X$;
the norm error in \eqref{eq:norm-error} is $C_{\rm app}X/\sqrt2$.
Either norm bound $O(X^{1+\delta/2})$ therefore implies the other
for $\delta\ge0$.
\end{proof}

The kernel $W_X$ is positive as a quadratic form, because it is a Gram
kernel. Its individual cross terms may have either sign. The exact
expression \eqref{eq:covariance} is the principal open arithmetic
formulation: bounding its terms individually by absolute values can
discard the cancellation one needs. Equation~\eqref{eq:Q-equivalence}
is an equivalence of admissible growth exponents, not an equality of
leading coefficients.

\begin{corollary}[Exponent budget for a proposed input]\label{cor:gate}
Let $h=X^\tau$, with fixed $0<\tau<1$. If an unconditional estimate
\begin{equation}\label{eq:candidate-S}
S(X,h)\le C Xh^2X^{-\kappa}L(X),\qquad \kappa>0,\quad
L(X)=X^{o(1)},
\end{equation}
holds for all sufficiently large real $X$, then
\begin{equation}\label{eq:power-budget}
\Vcal(X)\ll_g X^{3-\kappa}L(X)+X^{4\tau-1}.
\end{equation}
It follows that every $\delta>\max(0,1-\kappa,4\tau-3)$ is admissible.
\end{corollary}
\begin{proof}
Substitute into \eqref{eq:gate}, using $L(X)\ll_\epsilon X^\epsilon$
for every $\epsilon>0$. The displayed strict inequality allows
$\epsilon$ to be chosen smaller than the available exponent margin.
\end{proof}

Table~\ref{tab:budget} lists implications of hypothetical inputs;
none of its positive $\kappa$ estimates is proved here. At
$\tau=3/4$, all support and averaging costs are already $O(X^2)$.
An unbounded $L$ does not in general justify the frontier
$\delta=1-\kappa$. With $L=O(1)$, one can instead take the non-strict
maximum in the corollary. If only the averaging term determines the
frontier, its endpoint can also follow directly from
\eqref{eq:power-budget}.
\begin{table}[ht]
\centering
\begin{tabular}{cccc}
\toprule
Length exponent $\tau$ & Relative saving $\kappa$ & Sufficient $\delta$ & Consequence\\
\midrule
$3/4$ & $1/100$ & $\delta>0.99$ & A fixed strip if $\delta<1$\\
$3/4$ & $1/2$ & $\delta>1/2$ & Narrower fixed strips\\
$3/4$ & $1$ & Every $\delta>0$ & RH\\
$7/8$ & $1$ & $\delta>1/2$ & Averaging cost remains\\
\bottomrule
\end{tabular}
\caption{Translation of proposed raw mean-square savings with $L=X^{o(1)}$.}
\label{tab:budget}
\end{table}

\subsection{The available short-interval input retains exponent one}
The unconditional input audited from Saffari--Vaughan \cite[Section 6]{SV}
gives, at $h=X^{3/4}$,
\begin{equation}\label{eq:sv}
S(X,h)\ll_g Xh^2\exp[-c_g(\log X/\log\log X)^{1/3}].
\end{equation}
There is a distinction in the source: the stated additive Lemma 6 is
for $\theta$, whereas the proof first establishes a relative-interval
estimate for $\psi$ in (6.12)--(6.18). We use that latter estimate,
not a silent replacement of $\theta$ by $\psi$ in the lemma.

Here is the short transfer needed to specify the input's scope. Write
$I(u,v)=\psi(u+v)-\psi(u)-v$. The source's relative estimate, with
its unconditional density exponent $12/5$ and a fixed positive margin
in the length condition, gives
\[
\int_U^{2U}|I(u,\vartheta u)|^2\dd u
\ll\vartheta^2U^3 r(U),\qquad
r(U)=\exp[-c(\log U/\log\log U)^{1/3}],
\]
uniformly for $U\ge4$ and $U^{\epsilon-5/6}<\vartheta<1$, for
each fixed $\epsilon>0$; below take $\epsilon=1/12$.
For $2h\le v\le3h$,
$I(u,h)=I(u,v)-I(u+h,v-h)$. Squaring, integrating over
$u\in[U,2U]$ and $v\in[2h,3h]$, and assuming $h\le U/6$ gives
\[
h\int_U^{2U}|I(u,h)|^2\dd u
\le12U\int_{h/(3U)}^{3h/U}\int_U^{3U}
|I(u,\vartheta u)|^2\dd u\dd\vartheta.
\]
Cover $[U,3U]$ by $[U,2U]$ and $[2U,4U]$, apply the relative
estimate, and integrate $\vartheta^2$. This proves
$\int_U^{2U}|I(u,h)|^2\dd u\ll Uh^2r(U)$.
For $h=X^{3/4}$ and $U$ a fixed multiple of $X$, the entire
$\vartheta$ band is admissible: its lower end is comparable to
$X^{-1/4}$, above the source threshold $U^{-3/4}$ obtained by taking
the length margin $1/12$ with density exponent $12/5$.

Finally, $u=t-h/2$ moves the band of $S$ into $[AX/2,2BX]$.
The three dyadic intervals with lower endpoints
$U_j=2^jAX/2$, $j=0,1,2$, cover that band, since $4AX>2BX$.
This proves \eqref{eq:sv} with both ends included. Through
\eqref{eq:gate} it gives only
\[
\Vcal(X)\ll_g X^3\exp[-c_g(\log X/\log\log X)^{1/3}]+X^2.
\]
The saving is subpower and leaves $\delta=1$. This tested input does
not improve \eqref{eq:jy}; it does not rule out a stronger signed
covariance estimate.

\section{The first global milestone and the absence of a descent rule}
\label{sec:milestone}
A concrete first target is a fixed $\kappa>0$ in
\eqref{eq:candidate-S} at $h=X^{3/4}$, or a direct bound
$Q_{X^{3/4}}(X)=O(X^{2+\delta})$ for some fixed $\delta<1$.
All these estimates must concern the actual coefficients and all
sufficiently large real $X$. The support and approximation analysis
already has the required quadratic cost; the missing step is arithmetic.
An alternative useful target has the form
\[
\Vcal(X)\le C\{X^2L(X)+X^{3-\kappa}L(X)\},\qquad L=X^{o(1)}.
\]
It gives every $\delta>\max(0,1-\kappa)$, and would imply RH if
$\kappa\ge1$. This is a candidate bound, not an established estimate.

The following proposition explains the limits of using only generic
counting and preparation properties. It is a statement about constructed
coefficients, with no claim that they are Euler-prime coefficients.

\begin{proposition}[Finite-prefix obstruction]\label{prop:model}
Fix $0<\delta<1$, $\gamma\in\R\setminus\{0\}$, $0<\eta<1/2$, and an integer
$N_0\ge2$. There are nonnegative coefficients $c_n$, equal to
$\Lambda(n)$ for $2\le n\le N_0$ and belonging to $\{0,\log n\}$
beyond a finite index, such that, with $\beta=(1+\delta)/2$,
\begin{equation}\label{eq:model-psi}
\widetilde\psi(x):=\sum_{2\le n\le x}c_n
=x+\eta\Rea\frac{x^{\beta+i\gamma}}{\beta+i\gamma}+O(\log x).
\end{equation}
Define $\widetilde V(x)=\sum c_nw(n/x)$ and
$\widetilde{\mathcal V}(X)=\int_X^{2X}|\widetilde V(x)|^2\dd x$.
Then, for all sufficiently large real $X$,
\begin{equation}\label{eq:model-var}
\widetilde{\mathcal V}(X)\asymp X^{2+\delta}.
\end{equation}
The associated complete diagonal
\begin{equation}\label{eq:model-diagonal}
\widetilde D(Y)=\int_0^Y\sum_{n\ge2}\frac{c_n^2}{n}
g(y-\log n)^2\dd y
\end{equation}
satisfies $\widetilde D(Y)=Y^2/2+O(1)$.
\end{proposition}
The proof, including the finite-prefix and all-real-$X$ assertions, is
in Appendix~\ref{app:model}. These sequences satisfy PNT and even
$\widetilde\psi(x)-x=O(xe^{-cL(x)})$ for each fixed $c>0$ and
nonnegative $L=o(\log x)$. The latter is an asymptotic assertion with
parameter-dependent constants and thresholds, not preservation of every
published explicit constant. Their probe preparation is identical and
their diagonal has the same leading order, but they defeat every exponent
smaller than $\delta$. By choosing $N_0$ large enough they also preserve
any prescribed finite collection of tests depending only on a bounded
coefficient prefix.

Thus these controls and an input exponent alone cannot imply a universal
descent rule. Even a proposed universal step from $\delta=1$ to a fixed
$\delta'<1$ is refuted within this generic class by choosing
$\delta\in(\max(0,\delta'),1)$. This says nothing against a theorem
using the prime-power support, multiplicative structure, or additional
signed arithmetic information of the actual coefficients. The global
milestone remains open; the fixed-exponent theorem, finite certificates,
and transfer identities specify its exact requirements.

\appendix
\section{Certificate constants and reproducibility}\label{app:certificate}
The finite theorem uses a small transfer of existing outward records.
The following details make its analytic constants explicit without
duplicating a zero cache.

\subsection{Exact polynomial arithmetic}
Write $b_j=(-16)^j\binom8j$. On $(-1/4,1/4)$,
\[
g_0(v)=-\sum_{j=2}^8b_j(2j)(2j-1)(2j-2)v^{2j-3}
+\sum_{j=1}^8\frac{j b_j}{2}v^{2j-1}.
\]
For any polynomial product one integrates each monomial using
\[
\int_{-1/4}^{1/4}v^m\dd v=
\begin{cases}2\,4^{-(m+1)}/(m+1),&m\text{ even},\\0,&m\text{ odd}.
\end{cases}
\]
Squaring the displayed polynomial and its sixth derivative yields exactly
\[
\norm{g_0}_2^2=\frac{146640624550936576}{37921101075},\qquad
\norm{g_0^{(6)}}_{L^2(-a,a)}^2
=\frac{2504085215525628254072340480}{19}.
\]
The total variation of the two endpoint atoms of $D^6g_0$ is
$2\cdot8!\,8^8=1352914698240$. Cauchy--Schwarz bounds its interior
$L^1$ mass by the square root of half the last squared norm, which is
strictly less than $8117695446119$ (an integer-square comparison).
Consequently
\[
\norm{D^6g_0}_{\rm TV}<9470610144359=:M_6,\qquad K=M_6/\sqrt\nu
\]
is valid in Lemma~\ref{lem:decay}. These finite rational calculations
are independently reproducible from the polynomial prescription above.

\subsection{Zero counts and outward inputs}
The explicit count in \cite[Corollary 1.2]{HSW} is, for $T\ge e$,
\[
\left|N(T)-\frac{T}{2\pi}\log\frac{T}{2\pi e}\right|
\le0.1038\log T+0.2573\log\log T+9.3675.
\]
For $T\ge100$, $\log\log T\le\log T$ and
$9.3675\le(9.3675/\log100)\log T$.
The main term is at most $(T/(2\pi))\log T$.
The remaining coefficient is less than $2.4$, whereas
$(1-1/(2\pi))T>84$ at $T=100$ and increases. This verifies
$N(T)\le T\log T$ on the full range used in \eqref{eq:tail}.

The outward package \cite{OutwardCertificate} reconstructs the probe
and uses Arb/Acb zero enumeration and counting. Its 192-bit and 256-bit
records each check $N(500)=269$, enclose 270 consecutive positive zeros,
verify $\gamma_{269}<500<\gamma_{270}$ and ordered disjoint enclosures,
and bound the low mass including both signs. The count and consecutive
enumeration, not just a list of decimal approximations, provide
completeness; the routines are documented in \cite{FLINT}.
The published finite-height theorem \cite{PT} separately supplies
criticality through $H=3\cdot10^{12}$.

The finite transfer \cite{FiniteCertificate} checks the input hashes,
their scope, overlapping enclosures, the strict roundings
\eqref{eq:rounding}, and the following exact budgets:
\begin{center}
\begin{tabular}{ccc}
\toprule
Upper endpoint & Rational upper bound for $\Vcal(X)/X^2$ & Strict ceiling\\
\midrule
$10^{99}$ & $2955558057/78125000$ & $38$\\
$10^{100}$ & $62255877/1562500$ & $40$\\
$10^{102}$ & $5580201/78125$ & $72$\\
\bottomrule
\end{tabular}
\end{center}
For the cross coefficient it uses $(283/50)^2>32$; at $10^{99}$ it
also uses $\sqrt{10}<16/5$. All comparisons are rational.
The original records use the minus-exponent transform $G(-s)$.
Their coefficient masses are absolute values, so the convention change
does not alter the imported bounds.

Paths in the following list are relative to the parent directory
\path{../investigations/sonin-critical-boundary/numerics/}.
These SHA-256 identifiers locate the exact imported inputs and transfer:
\begin{description}
\item[Outward generator] \path{selective_loss_quadratic_target_20261003/certify.py}\\
\nolinkurl{6f3947f99d4b45307b8a84fb25a6ec56d395d98fb5467106d77dac0a95c63a1b}
\item[192-bit record] \path{selective_loss_quadratic_target_20261003/bound_192.json}\\
\nolinkurl{21a71598fb9dcf8fe51743e75f65aedd5c5deee69619d2bda806da3e9ac6f2c0}
\item[256-bit record] \path{selective_loss_quadratic_target_20261003/bound_256.json}\\
\nolinkurl{3c39ccfbc704122de57c04395e120f11a5a563d9d3838c5b6388a9e7c50a60b8}
\item[Finite transfer generator] \path{subpower_finite_variance_20261003/certify.py}\\
\nolinkurl{91e1b31ce6a34d6f6625c1b6ccb2c02e60c4c404599277ba274bfd9e3eafbf1f}
\item[Finite transfer record] \path{subpower_finite_variance_20261003/record.json}\\
\nolinkurl{7576708c6cd4e0e9be89d7551d1386ef77b9c972908f55d50f864b9c4afc6b26}
\end{description}
The small generator was replayed into a temporary output and matched the
retained finite record byte-for-byte; the exact polynomial arithmetic
was also checked. This replay imports the earlier outward calculations;
it does not independently regenerate them or the published verification.
Hashes identify data, while the rational inequalities and the analytic
proof establish the transfer. The linked package READMEs give the
reproduction commands and software requirements.

\section{Proof of the finite-prefix obstruction}\label{app:model}
\begin{proof}[Proof of Proposition~\ref{prop:model}]
Put $N=\max(N_0,8)$ and retain $c_n=\Lambda(n)$ through $N$.
Let $S_N=\sum_{2\le n\le N}\Lambda(n)$ and set
\[
M(x)=x+\eta\Rea\frac{x^{\beta+i\gamma}}{\beta+i\gamma},\qquad
F(x)=S_N+M(x)-M(N)\quad(x\ge N).
\]
Since $\beta<1$, $1-\eta\le M'(x)\le1+\eta<2$ for $x\ge1$.
For each integer $n>N$, choose $c_n=\log n$ if $S_{n-1}<F(n)$,
and $c_n=0$ otherwise; put $S_n=S_{n-1}+c_n$.
Induction proves
\begin{equation}\label{eq:tracking}
0\le S_n-F(n)\le\log n\qquad(n\ge N).
\end{equation}
Indeed the increment $d_n=F(n)-F(n-1)$ lies in $(0,2)$. If the previous
excess minus $d_n$ is nonnegative, no insertion is made and the bound
persists. Otherwise insertion gives excess in $(\log n-2,\log n)$,
which is positive because $n>N\ge8$. Between successive integers
$F$ changes by less than two. Thus
$\widetilde\psi(x)=F(x)+O(\log x)=M(x)+O(\log x)$, proving
\eqref{eq:model-psi} including the prescribed prefix.

For the complete support of $w(t/x)$, Stieltjes integration against
$dM$ cancels its constant-density part. If $s=\beta+i\gamma$,
substitution $u=e^{-v}$ gives the exact Mellin identity
\[
\int_A^B w(u)u^{s-1}\dd u=G(1/2-s).
\]
The tracking error contributes $O(\log x)$ by integration by parts:
it is bounded by
$x^{-1}\int_{Ax}^{Bx}O(\log t)|w'(t/x)|\dd t$.
Endpoint terms vanish, including for the full finite windows.
Writing $\alpha=\delta/2$, we obtain
\begin{equation}\label{eq:model-response}
\widetilde V(x)=\eta\Rea\{G(-\alpha-i\gamma)x^{\beta+i\gamma}\}
+O(\log x).
\end{equation}
By Lemma~\ref{lem:noncancel} its complex amplitude is nonzero.
Let $\lambda=2+\delta$, $\ell=\log2$, and
$\theta_X=\gamma\log X+\arg G(-\alpha-i\gamma)$.
Using $x=Xe^t$ in the squared main term gives
\begin{equation}\label{eq:phase}
\begin{split}
\widetilde{\mathcal V}(X)
&=\eta^2|G(-\alpha-i\gamma)|^2X^\lambda Q(\theta_X)\\
&\quad+O\bigl(X^{3/2+\delta/2}\log X+X(\log X)^2\bigr),\\
Q(\theta)&=\int_0^\ell e^{\lambda t}\cos^2(\theta+\gamma t)\dd t.
\end{split}
\end{equation}
Both error powers are smaller than $\lambda$. Direct integration yields
\begin{equation}\label{eq:phase-min}
\min_\theta Q(\theta)=\frac12\left\{
\frac{e^{\lambda\ell}-1}{\lambda}
-\left|\frac{e^{(\lambda+2i\gamma)\ell}-1}{\lambda+2i\gamma}\right|
\right\}>0.
\end{equation}
The strict inequality is the strict triangle inequality for
$\int_0^\ell e^{\lambda t}e^{2i\gamma t}\dd t$: its phase is not
constant on any interval because $\gamma\ne0$.
The maximum of $Q$ is finite. The uniform positive minimum in
\eqref{eq:phase-min}, not merely a favorable subsequence of phases,
proves \eqref{eq:model-var} for every sufficiently large real $X$.

Finally put $\widetilde A_2(t)=\sum_{\log n\le t}c_n^2/n$.
Since $c_n^2=(\log n)c_n$ for $n>N$, there is a fixed prefix correction
\[
C_N=\sum_{2\le n\le N}\frac{c_n^2-(\log n)c_n}{n}
\]
such that, once the prefix is fully included,
\begin{equation}\label{eq:diagonal-prefix}
\widetilde A_2(t)=\int_{2^-}^{e^t}\frac{\log u}{u}\dd\widetilde\psi(u)
+C_N,\qquad t\ge\log N.
\end{equation}
Partial summation with
$\widetilde\psi(u)-u=O(u^\beta)+O(\log u)$ gives
$\widetilde A_2(t)=t^2/2+C+o(1)$: the endpoint error
$O(t e^{-(1-\beta)t})$ tends to zero and the differentiated-error
integral converges because $\beta<1$.
Every packet begins above zero since $a<\log2$. Finite interchange
therefore proves the exact identity
\[
\widetilde D(Y)=\int_{-a}^a g(v)^2\widetilde A_2(Y-v)\dd v.
\]
For large $Y$ all arguments exceed $\log N$.
Use $\int g^2=1$ and $\int v g(v)^2\dd v=0$ to conclude
$\widetilde D(Y)=Y^2/2+O(1)$. The retained prefix affects only the
constant. This completes the proof.
\end{proof}

\section*{Acknowledgements and status}
This draft was prepared with substantial large-language-model assistance
in mathematical development, exposition, source checking, and internal
proof review. Internal checks do not constitute independent specialist
refereeing. The analytic reductions are proved in the text; the finite
numerical theorem additionally relies on the published verification and
the identified outward records. No global exponent below one and no
actual-prime exponent-descent mechanism is claimed.

\begin{thebibliography}{99}
\bibitem{DLMF}
F.~W.~J.~Olver et al., eds.,
\emph{NIST Digital Library of Mathematical Functions},
Sections 25.2, 25.4, 25.10, and equation 27.4.12.
\url{https://dlmf.nist.gov/25.4};
\url{https://dlmf.nist.gov/27.4.E12}.

\bibitem{ExplicitFormula}
T.~Tao, \emph{254A, Notes 2: Complex-analytic multiplicative number
theory}, 2014, especially Propositions 16 and 19 and the discussion
of explicit formulae.
\url{https://terrytao.wordpress.com/2014/12/09/254a-notes-2-complex-analytic-multiplicative-number-theory/}.

\bibitem{JY}
D.~R.~Johnston and A.~Yang,
\emph{Some explicit estimates for the error term in the prime number theorem},
J. Math. Anal. Appl. \textbf{527} (2023), no.~2, Article 127460.
\href{https://doi.org/10.1016/j.jmaa.2023.127460}{doi:10.1016/j.jmaa.2023.127460};
\href{https://arxiv.org/abs/2204.01980v2}{author preprint}, Theorem 1.4.

\bibitem{PT}
D.~J.~Platt and T.~S.~Trudgian,
\emph{The Riemann hypothesis is true up to $3\cdot10^{12}$},
Bull. London Math. Soc. \textbf{53} (2021), no.~3, 792--797.
\href{https://doi.org/10.1112/blms.12460}{doi:10.1112/blms.12460};
\href{https://arxiv.org/abs/2004.09765v1}{author preprint}, Theorem 1.

\bibitem{HSW}
E.~Hasanalizade, Q.~Shen, and P.-J.~Wong,
\emph{Counting zeros of the Riemann zeta function},
J. Number Theory \textbf{235} (2022), 219--241.
\href{https://doi.org/10.1016/j.jnt.2021.06.032}{doi:10.1016/j.jnt.2021.06.032};
\href{https://arxiv.org/abs/2107.06506v1}{author preprint}, Corollary 1.2.

\bibitem{SV}
B.~Saffari and R.~C.~Vaughan,
\emph{On the fractional parts of $x/n$ and related sequences. II},
Ann. Inst. Fourier \textbf{27} (1977), no.~2, 1--30, Section 6.
\href{https://doi.org/10.5802/aif.649}{doi:10.5802/aif.649}.

\bibitem{FLINT}
The FLINT contributors, \emph{FLINT documentation: acb\_dirichlet},
zero-counting and consecutive zeta-zero enclosure routines.
\url{https://flintlib.org/doc/acb_dirichlet.html}.

\bibitem{OutwardCertificate}
\emph{Quadratic signed-pair target: outward numerical certificate},
project numerical package, 3 October 2026.
\href{../investigations/sonin-critical-boundary/numerics/selective_loss_quadratic_target_20261003/README.md}{Package README};
\path{../investigations/sonin-critical-boundary/numerics/selective_loss_quadratic_target_20261003/README.md}.
Includes the original 192-bit and 256-bit records and their generator.

\bibitem{FiniteCertificate}
\emph{Finite continuous dyadic variance certificate},
project numerical package, 3 October 2026.
\href{../investigations/sonin-critical-boundary/numerics/subpower_finite_variance_20261003/README.md}{Package README};
\path{../investigations/sonin-critical-boundary/numerics/subpower_finite_variance_20261003/README.md}.
Includes the small rational transfer and its exact record.
\end{thebibliography}
\end{document}
